\clearpage
\section{Literature Review}
\label{sec:literature}
\subsection{Human Capital and Segmented Labor Market Theories}

Returns to education and human capital have been subject to extensive empirical and theoretical investigation in economics. One of the first papers to formalise the interdependence between education and individual earnings was \shortciteA{Mincer1958}, who later established that the logarithm of individual earnings can be expressed as a linear function of years of schooling and a quadratic in experience \cite{Mincer1974}. This earnings function became the dominant empirical framework in labour economics for half a century, yet it rests on strong assumptions---constant marginal returns to each year of schooling, separability of education and experience---that subsequent work has systematically questioned \cite{HeckmanLochnerTodd2006}.

Becker's (\citeyear{Becker1964}) theoretical framework predicts that the rate of return to college education should vary systematically across demographic groups once individual differences in financing costs and opportunity costs are accounted for. Empirically, he documented that estimated returns to schooling are positive even for historically disadvantaged groups, although the raw college wage premium was larger for white male cohorts---a pattern he attributed to unobserved heterogeneity in school quality and credit constraints rather than differences in underlying productive capacity \cite[Ch.~4]{Becker1964}. This founding insight motivates the present study: if structural barriers rather than human capital deficits drive differential returns, then access-expanding policies like ProUni face an institutional ceiling that increases in education supply alone cannot lift.

In a comprehensive survey and meta-analysis of the causal returns to education literature, \citeA{Card1999} confirmed the positive direction of these results using causal inference methods and found that each additional year of schooling raises earnings by approximately 7--10\%. Crucially, \citeA{Card1999} highlights that policy-based instruments---school construction, proximity to colleges, financial aid rules---typically yield Local Average Treatment Effect (LATE) estimates around 20\% above OLS, reflecting the particularly high marginal returns earned by students who change their educational choices in response to such policies. He also highlights that ``sheepskin'' effects \cite{Hungerford1987}---the wage discontinuity at credential completion---may compound productivity gains with signalling premiums, and that these effects interact with ethnicity and gender \cite{BelmanHeywood1991}.
\paragraph{}
While \emph{human capital} theory interprets wage differentials as the result of individual investment outcomes in competitive markets, \emph{segmented labour market} theory complements this view by emphasising institutional factors, hierarchies, and unequal access to the primary labour market \cite{ReichGordonEdwards1973, DoeringerPiore1971, Leontaridi1998}. The theory conceptualised in \citeA{ReichGordonEdwards1973} argues that the labour market is structurally divided into a primary sector---characterised by high wages, stability, and internal promotion ladders---and a secondary sector---characterised by low wages, instability, and limited returns to credentials. Crucially, the allocation of workers to sectors is not solely determined by productivity but also by race, gender, and geographic origin, making segmentation a mechanism of structural reproduction rather than a neutral market outcome.

\citeA{Leontaridi1998} offers a synthetic assessment of the segmented labour market literature: while the theory captures important institutional realities, its explanatory power is weakened by inconsistent empirical definitions of what constitutes a ``segment'' and by difficulty in separating segmentation from compensating differentials. She recommends integrating segmented labour market insights with human capital and search models to produce more tractable empirical predictions. This methodological caution applies directly to the present thesis: the fact that ProUni operates through the private sector---structurally positioned in the secondary tier of Brazil's dual university system---raises the question of whether its graduates gain access to the primary labour market or remain confined to the lower-return segment.
\paragraph{}
The conclusions of both theories make it important to study the institutional mechanisms that drive the apparent heterogeneity in educational returns. If segmentation drives return differentials along demographic lines, then education policies solely aimed at expanding access may have a limited effect on reducing inequality or lifting people out of poverty. I therefore expect \textbf{greater effects} where policies open access to selective programmes and dense professional networks, thereby lowering ethnic and gender gaps, and \textbf{smaller effects} where regional labour demand limits the conversion of credentials into high-quality employment. Since most of the foundational studies cited above were conducted in the U.S.~labour market, Section~\ref{subsec:international_evidence} surveys comparable evidence in international contexts, and Section~\ref{subsec:brazil_returns} examines the specific Brazilian evidence.




\subsection{Returns to Education in Brazil}
\label{subsec:brazil_returns}
The analysis of educational returns in Brazil must be contextualised within the country's profound structural duality, famously termed \textbf{\emph{Belíndia}} by the economist Edmar Bacha in the 1970s \cite[p.~19]{Telles2004}. This blended word highlights the coexistence of a ``small, growing industrial state such as Belgium'' alongside a ``great mass of poor stagnating at almost the level of rural India'' \cite{bacha1976}. The entrenched economic segmentation justifies why educational policies must simultaneously address individual human capital deficits and the deep structural barriers that restrict social mobility between these two ``countries.'' Rather than treating economic agents as a homogeneous entity, this framework naturally leads to an analysis of the specific ethnic, gender, and regional market failures observed in Brazil.

In the context described above, several contributions to Brazilian social science have employed methodologies in the tradition of \citeA{Mincer1974} to estimate educational returns with a pronounced focus on heterogeneity.
\paragraph{}
First, evidence shows ethnic wage gaps beyond observables and a skin colour gradient. \citeA[p.~171]{Telles2004} documents a degree of income-distribution disparity comparable only to South Africa among major developing economies. In his meta-study of racial discrimination, \citeA[pp.~156--163]{Telles2004} identifies a glass-ceiling mechanism separating white from non-white workers that almost entirely blocks the latter group from ascending into the middle class. \citeA{AriasYamadaTejerina2002} demonstrate an unexplained earnings gap of approximately 18\% between non-white and white university graduates, based on data covering 1940--1990, using historical pupil-teacher ratios as instruments for educational quality. Using matched employer--employee data from RAIS and PNAD, \citeA{Cornwell2017} confirm the persistence of this gap in the formal sector even after conditioning on observable worker and firm characteristics, providing direct evidence of labour-market discrimination beyond selection.
\paragraph{}
Second, there is evidence of persistent gender gaps even though women attain equal or higher levels of formal education. \citeA{BusteloEtAl2021} document that the Brazilian gender pay gap is substantially explained by occupational sorting at the point of college major selection: women disproportionately choose Health, Teaching, and Social Science fields, which carry lower average wages in the formal sector, while men sort into Engineering, Technology, and Business programmes with higher returns. \citeA{GarciaNopoSalardi2009} show using Ñopo's non-parametric decomposition that this sorting effect, combined with within-occupation discrimination captured by the unexplained component, sustained gender wage gaps of 20--25\% throughout the 1996--2006 period. These findings directly predict the asymmetric gender results observed in this thesis: a policy that increases access uniformly does not overcome the downstream sorting and structural barriers that determine how credentials are priced in the labour market \cite{Goldin2014}.
\paragraph{}
Third, there are sizable regional differences in labour demand and in the conversion of education into wages, documented by \citeA{deoliveira2021}. Average labour income in the metropolitan area of Fortaleza (Northeast) is more than 60\% below the national average of Brazilian metropolitan regions, and \citeA{deoliveira2021} demonstrate that this gap carries a positive gender dimension: women face a compounded disadvantage of lower regional wages and a larger unexplained within-region gender gap. This evidence reinforces the thesis's hypothesis that the effects of educational expansion depend critically on local labour-demand conditions---the ``Brazil'' side of Belíndia may be precisely where marginal returns to skilled labour are highest, yet where private HEI infrastructure is thinnest.
\paragraph{}
Policies can reshape these margins \emph{unevenly}. \citeA{RiehlMachadoReyes2023} provide the most direct Brazilian evidence using a two-stage regression discontinuity design to study the direct and spillover effects of the 45\% affirmative action quota at UERJ, an elite public university in Rio de Janeiro. They find a 14\% increase in the early-career earnings of marginally admitted non-white beneficiaries, attributable primarily to network access to high-paying firms rather than to human capital acquisition per se. Critically, however, these gains are decreasing in career progression---consistent with structural glass-ceiling barriers reasserting themselves as graduates age---and the spillover effect on top-ranked non-beneficiary enrollees in high-take-up majors is negative: early-career earnings for that group fall by 14\%. Their results also confirm a spatial gradient in programme returns: gains are concentrated in the dense, productive labour market of São Paulo, while graduates in the Northeast capture substantially smaller premia. This spatial heterogeneity is a central motivation for the microregion-level design in this thesis.



\subsection{International Evidence on Education Policies}
\label{subsec:international_evidence}

The international evidence on large-scale education supply shocks offers two complementary lessons for the ProUni context. On one hand, Duflo's (\citeyear{Duflo2001}) analysis of Indonesia's INPRES school construction programme---the closest institutional analogue to ProUni in terms of identification strategy---shows that regions receiving greater construction intensity experienced larger increases in schooling years and wages for exposed birth cohorts, with estimated annual returns of up to 10\%. Crucially, \citeA{Duflo2001} identifies these returns through the interaction of \emph{spatial intensity} of the infrastructure shock with \emph{cohort exposure}---the same two-dimensional identification logic employed in this thesis. A key finding is that the returns are concentrated among cohorts young enough to actually attend the new schools; older cohorts in the same districts show no effect, a pattern that motivates the age-cohort placebo design in Chapter~\ref{sec:methodology}.

A broader meta-analysis by \citeA{PsacharopoulosPatrinos2004} across 139 countries confirms that returns to tertiary education are generally positive and often highest in developing economies where skilled labour is relatively scarce---consistent with the scarcity-versus-saturation hypothesis developed in this thesis.

Evidence from selective institutional settings reveals important distributional limits. \citeA{BagdeEppleTaylor2016} show that India's caste-based reserved seats altered college quality access and occupational sorting for socially disadvantaged groups, but that institutional quality acted as a binding constraint on the realised returns: beneficiaries assigned to lower-quality institutions captured smaller wage gains than those assigned to elite colleges. In the U.S. context, \citeA{Bleemer2022} finds that the elimination of affirmative action in California reduced minority enrolment at selective universities but also diminished downstream earnings, confirming that institutional selectivity is a key mediator of policy returns. \citeA{ArcidiaconoLovenheim2020} synthesise this evidence and argue that the quality-fit trade-off---whether beneficiaries are better served by a less selective institution where they perform well, or a more selective one where they struggle---depends critically on the local labour market and employer preferences.

Taken together, these studies establish three expectations that guide the empirical design of this thesis: (i) supply-side expansions generate positive wage returns for exposed cohorts when local labour demand is sufficient to absorb graduates; (ii) the magnitude of returns is highly heterogeneous across regions, demographic groups, and institutional quality tiers; and (iii) large-scale programmes may generate general equilibrium effects---including congestion and spillovers---that attenuate private returns in heavily exposed markets. The following chapters test each of these predictions in the Brazilian context.
